%%%%%%%%%%%%%%%%%%%%%%%%%%%%%%%%%%%%%%%%%%%%%%%%%%%%%%%%%%%%%%%%%%%%%%%%%%%%%%
%% A five-condition uniqueness theorem for the gauge sector of the Standard Model
%% Author: Kevin Remondiere
%% Target: Letters in Mathematical Physics (LMP), alt. Annales Henri Poincare
%% License: CC-BY 4.0
%% Date: 2026-05-24
%%%%%%%%%%%%%%%%%%%%%%%%%%%%%%%%%%%%%%%%%%%%%%%%%%%%%%%%%%%%%%%%%%%%%%%%%%%%%%

\documentclass[11pt,reqno]{amsart}

\usepackage[utf8]{inputenc}
\usepackage[T1]{fontenc}
\usepackage[english]{babel}
\usepackage{amsmath,amssymb,amsthm,mathtools}
\usepackage{geometry}
\geometry{a4paper,margin=1in}
\usepackage{hyperref}
\hypersetup{
  colorlinks=true,
  linkcolor=blue,
  citecolor=blue,
  urlcolor=blue,
}
\usepackage{enumitem}
\usepackage{microtype}
\usepackage{booktabs}
\usepackage{array}

%% Theorem environments
\theoremstyle{plain}
\newtheorem{theorem}{Theorem}[section]
\newtheorem{proposition}[theorem]{Proposition}
\newtheorem{lemma}[theorem]{Lemma}
\newtheorem{corollary}[theorem]{Corollary}

\theoremstyle{definition}
\newtheorem{definition}[theorem]{Definition}
\newtheorem{condition}[theorem]{Condition}

\theoremstyle{remark}
\newtheorem{remark}[theorem]{Remark}

%% Custom macros
\newcommand{\Tr}{\operatorname{Tr}}
\newcommand{\tr}{\operatorname{tr}}
\renewcommand{\Re}{\operatorname{Re}}
\newcommand{\dd}{\mathrm{d}}
\newcommand{\bbR}{\mathbb{R}}
\newcommand{\bbZ}{\mathbb{Z}}
\newcommand{\bbN}{\mathbb{N}}
\newcommand{\bbP}{\mathbb{P}}
\newcommand{\calL}{\mathcal{L}}
\newcommand{\calE}{\mathcal{E}}
\newcommand{\calP}{\mathcal{P}}
\newcommand{\KL}{D_{\mathrm{KL}}}
\newcommand{\TV}{\mathrm{TV}}
\newcommand{\Lat}{\Lambda}
\newcommand{\rk}{\operatorname{rk}}
\newcommand{\sun}{\mathfrak{su}(N)}
\newcommand{\SUN}{\mathrm{SU}(N)}
\newcommand{\SU}{\mathrm{SU}}
\newcommand{\msf}[1]{\mathsf{#1}}

\newcommand{\Cone}{(\mathrm{C}_1)}
\newcommand{\Ctwo}{(\mathrm{C}_2)}
\newcommand{\Cthree}{(\mathrm{C}_3)}
\newcommand{\Cfour}{(\mathrm{C}_4)}
\newcommand{\Cfive}{(\mathrm{C}_5)}

\title[Five-condition uniqueness for SM gauge sector]{A five-condition uniqueness theorem for the gauge sector of the Standard Model}

\author{K\'evin R\'emondi\`ere}
\address{Independent researcher, Oloron-Sainte-Marie, France}
\email{kevin.remondiere@gmail.com}
\thanks{ORCID: 0009-0008-2443-7166. Manuscript v1, 2026-05-24. Distributed under CC-BY 4.0.}

\date{\today}

\subjclass[2020]{Primary 22E70; Secondary 81T13, 81R40, 81V22}

\keywords{Standard Model, gauge group, Lie group classification, Wilson lattice, log-Sobolev inequality, Hodge decomposition, Bott periodicity, anthropic principle}

\begin{document}

\begin{abstract}
We formulate five structural conditions $\Cone$--$\Cfive$ on pairs $(G, D)$ consisting of a compact connected Lie group $G$ and a spacetime dimension $D \in \bbZ_{\ge 1}$, and we prove that exactly one pair satisfies the joint conditions $\Cone$--$\Cfour$ among all simple compact Lie groups of rank $\le 8$ and their semi-simple products of total rank $\le 2$: the pair $(\SU(3), 4)$. Condition $\Cfive$ then forces the augmentation of $\SU(3)$ by a non-saturated sub-sector, of which the simplest realisation is the Standard Model electroweak content $\SU(2)_L \times \mathrm{U}(1)_Y$. The five conditions are: $\Cone$ logarithmic Sobolev saturation of the Wilson Gibbs measure; $\Ctwo$ coincidence of Lie-algebraic and Hodge-geometric realisations of the saturation coefficient $\kappa$; $\Cthree$ non-triviality of the dynamics ($D \ge 3$); $\Cfour$ admissibility of complex Weyl chirality via Bott periodicity ($D \equiv 4 \pmod 8$); $\Cfive$ existence of a non-saturated sub-sector enabling spontaneous symmetry breaking and Sakharov baryogenesis. The proof is by exhaustive enumeration over the Cartan--Killing classification, combined with a standard tensorisation argument for products. The theorem is structural: it does not predict masses, Yukawa couplings, or the chirality $V{-}A$ versus $V{+}A$; it does fix the gauge content of any internally consistent four-dimensional gauge theory with the saturation+chirality+SSB scaffolding.
\end{abstract}

\maketitle

\setcounter{tocdepth}{1}
\tableofcontents

%%%%%%%%%%%%%%%%%%%%%%%%%%%%%%%%%%%%%%%%%%%%%%%%%%%%%%%%%%%%%%%%%%%%%%%%%%%%%%
\section{Introduction}
\label{sec:intro}

\subsection{Context.}
The Standard Model (SM) of particle physics has the gauge content
\[
  G_{\mathrm{SM}} \;=\; \SU(3)_c \times \SU(2)_L \times \mathrm{U}(1)_Y,
\]
acting on a four-dimensional spacetime. To date this content is \emph{observed} and \emph{postulated}, not \emph{derived} from a first principle. Successive attempts at ``embedding into a larger group'' (Pati--Salam, Georgi--Glashow $\SU(5)$, Fritzsch--Minkowski $\mathrm{SO}(10)$, $E_6$, $E_8$) have explained part of the matter content under a unification assumption, but have left untouched the question of why the gauge factor is $\SU(3)_c \times \SU(2)_L \times \mathrm{U}(1)_Y$ rather than any other combination of compact Lie groups, and why spacetime is four-dimensional.

\subsection{The result.}
We identify five structural conditions $\Cone$--$\Cfive$ on pairs $(G, D)$ and prove:

\begin{theorem}[Main Theorem; restated as Theorem~\ref{thm:main} below]
Among all simple compact connected Lie groups $G$ of rank $\le 8$, and all semi-simple products of such groups of total rank $\le 2$, paired with $D \in \{1, 2, 3, 4, 5, 6, 7, 8\}$, the unique pair satisfying $\Cone$ through $\Cfour$ jointly is $(G, D) = (\SU(3), 4)$. Condition $\Cfive$ is automatically satisfied if and only if $G$ is augmented by a non-saturated sub-sector of total rank $\ge 1$.
\end{theorem}

The simplest realisation of $\Cfive$ that introduces only rank-$1$ generators with chirality structure is $\SU(2)_L \times \mathrm{U}(1)_Y$, the SM electroweak sector. Larger non-saturated augmentations (e.g.\ an additional $\mathrm{U}(1)_{B-L}$, dark photon, hidden sector) are not excluded by our theorem.

\subsection{Honest scope.}
This paper does \emph{not} claim to derive the SM from first principles in any complete sense. We do \emph{not} derive: fermion content, generation count, Yukawa couplings, CKM matrix, neutrino mass scale, the sign of the $V{-}A$ chirality versus $V{+}A$, nor the numerical value of the baryon asymmetry $\eta_B$. We \emph{do} derive: the gauge content $\SU(3)_c$ together with the necessity of a rank-$1$ non-saturated electroweak augmentation, and the spacetime dimension $D = 4$, both jointly fixed by a structural argument whose five inputs are explicitly named and individually justified.

\subsection{Outline.}
\S\ref{sec:five-conditions} states the five conditions $\Cone$--$\Cfive$ and motivates each. \S\ref{sec:uniqueness} contains the uniqueness theorem and its proof. \S\ref{sec:empirical} surveys the empirical evidence for $\Cone$ and $\Ctwo$ from a recent Hamiltonian Monte Carlo (HMC) campaign on $\SU(3)$ Wilson lattices, with Lean~4 formal certification of the algebraic coefficient $\kappa = 1/6$. \S\ref{sec:consequences} collects four falsifiable numerical consequences. \S\ref{sec:scope} articulates the honest scope and lists three distinct falsification routes. \S\ref{sec:anthropic} places the result in the multiverse / anthropic landscape literature. \S\ref{sec:ack} acknowledges sources and discloses LLM-assisted adversarial review per COPE guidelines.

%%%%%%%%%%%%%%%%%%%%%%%%%%%%%%%%%%%%%%%%%%%%%%%%%%%%%%%%%%%%%%%%%%%%%%%%%%%%%%
\section{The five conditions}
\label{sec:five-conditions}

Throughout the paper, $G$ denotes a compact connected Lie group with Lie algebra $\mathfrak{g}$, rank $r := \rk(G)$, positive-root system $\Phi^+ \subset \mathfrak{g}^*$, and dual Coxeter number $h^\vee$. We write $D \in \bbZ_{\ge 1}$ for the spacetime dimension and reserve the symbol $d$ for the lattice dimension when it equals $D$.

\subsection{$\Cone$ Saturation}\label{ssec:C1}

\begin{condition}[Saturation]\label{cond:C1}
Let
\[
  \Sigma(D) \;:=\; \frac{D(D-1)(5-D)}{6}.
\]
The pair $(G, D)$ is \emph{saturated} when $\rk(G) = \Sigma(D)$.
\end{condition}

The polynomial $\Sigma(D)$ arises in the Wilson lattice formulation of $G$-gauge theory as the dimension of the \emph{normal cone} to the gauge orbit of the trivial connection in $G^{E(\Lat_a)}$, modulo the Faddeev--Popov ghost determinant; equivalently, as the codimension of the algebraic boundary of the saturated regime in which a logarithmic-Sobolev (LSI) bound is conjecturally sharp. Concretely:

\begin{center}
\begin{tabular}{cccp{8cm}}
\toprule
$D$ & $\Sigma(D)$ & Saturated $G$ (compact, simple) \\
\midrule
$1$ & $0$ & trivial (no gauge content) \\
$2$ & $1$ & $\SU(2) = A_1$ (rank $1$) \\
$3$ & $2$ & $\SU(3) = A_2$ (rank $2$) \\
$4$ & $2$ & $\SU(3) = A_2$ (rank $2$) \\
$5$ & $0$ & trivial \\
$6$ & $-5$ & no solution (negative) \\
$\ge 7$ & $\le -7$ & no solution \\
\bottomrule
\end{tabular}
\end{center}

The saturation regime $\Sigma(D) \ge 1$ is therefore confined to $D \in \{2, 3, 4\}$. For $D = 4$, the unique simple compact group of rank $2$ in the saturated regime is $A_2 = \SU(3)$ (the alternatives $B_2 = \mathrm{Sp}(2)$ and $G_2$ have rank $2$ but a strictly larger root system, see Lemma~\ref{lem:cartan-killing} below).

\paragraph{Operational consequence.}
A conditional theorem (in the spirit of \cite{PaperLMP_E}) establishes that under a named concentration axiom, the Langevin generator of $\mu_{a, L, \beta}$ on $G^{E(\Lat_a)}$ admits a spectral gap of the form
\[
  \lambda_1(\calL_\beta) \;\ge\; \varepsilon(N, D) \cdot (1 - \kappa(G, D)) \cdot \beta \cdot L^{-2},
\qquad
  \kappa(G, D) \;:=\; \frac{1}{2 |\Phi^+(G)|},
\]
for saturated pairs. This is the dynamical content of $\Cone$: the Wilson measure admits an LSI uniformly in $\beta$ inside the saturated family, with multiplicative Lie-algebraic correction $\kappa$.

\subsection{$\Ctwo$ Double degeneracy (Lie--Hodge collapse)}\label{ssec:C2}

The saturation coefficient $\kappa$ admits \emph{two} a priori distinct realisations:
\begin{itemize}
  \item[(alg.)] $\kappa_{\mathrm{Lie}}(G) = 1/(2 |\Phi^+(G)|)$;
  \item[(geo.)] $\kappa_{\mathrm{Hodge}}(D) = 1/(2(D-1))$ on a closed oriented $D$-manifold of signature zero.
\end{itemize}

\begin{condition}[Lie--Hodge collapse]\label{cond:C2}
The pair $(G, D)$ satisfies the collapse condition if and only if
\[
  |\Phi^+(G)| \;=\; D - 1.
\]
\end{condition}

For $G = \SU(3) = A_2$ one has $|\Phi^+| = 3$, and $\Ctwo$ becomes $D = 4$. The geometric side rests on Hodge self-duality of $\Lambda^2 T^*M$ on a closed oriented $4$-manifold of vanishing signature; the algebraic side rests on Peter--Weyl decomposition of $L^2(G)$ into matrix coefficients indexed by $\Phi^+$.

\paragraph{Operational consequence.}
When $\Cone \wedge \Ctwo$ holds, the value of $\kappa$ is doubly determined and agrees: $\kappa(G, D) = 1/(2(D-1)) = 1/(2|\Phi^+|)$. The numerical predictions of \S\ref{sec:consequences} depend only on this common value.

\subsection{$\Cthree$ Non-triviality}\label{ssec:C3}

\begin{condition}[Non-triviality]\label{cond:C3}
$D \ge 3$.
\end{condition}

Two-dimensional Yang--Mills theory is exactly solvable via the heat kernel on $G$ (Driver--Sengupta--L\'evy \cite{Driver1989, Sengupta1992, Levy2003}), has no continuous Higgs phase, and possesses no mass gap as a confining $4$d theory does. The non-trivial content of $\Cone$--$\Ctwo$ is therefore meaningful only for $D \ge 3$; in $D = 2$ the conditions are vacuous but the dynamics is degenerate.

\subsection{$\Cfour$ Complex Weyl chirality (Bott periodicity)}\label{ssec:C4}

The existence of complex Weyl (chiral) spinors on a $D$-dimensional pseudo-Riemannian manifold is governed by the Atiyah--Bott--Shapiro classification of Clifford modules \cite{ABS64}, see also Lawson--Michelsohn \cite{LawsonMichelsohn}, Ch.~I~\S5. Complex Weyl spinors split into chiral halves precisely when $D$ is even, but the chiral halves are \emph{inequivalent} (genuinely complex and not pseudoreal) only when $D \equiv 2 \pmod 8$ or $D \equiv 4 \pmod 8$; further, anomaly-free chirality in a \emph{Lorentzian} signature requires $D \equiv 4 \pmod 8$.

\begin{condition}[Complex Weyl admissibility]\label{cond:C4}
$D \equiv 4 \pmod 8$.
\end{condition}

\begin{remark}
$\Cfour$ alone selects $D \in \{4, 12, 20, \dots\}$. Combined with $\Cthree$ and the saturation regime $\Sigma(D) \ge 1$ of $\Cone$, the value $D = 4$ is uniquely fixed.
\end{remark}

\subsection{$\Cfive$ Asymmetry YM--EW}\label{ssec:C5}

A theory satisfying $\Cone \wedge \Ctwo$ saturates at every $G$-orbit, which forbids the spontaneous breaking of $G$ to a smaller subgroup. To accommodate the Higgs mechanism and Sakharov's three baryogenesis conditions \cite{Sakharov1967} we require the existence of a sub-sector $G' \subseteq G_{\mathrm{total}}$ that is \emph{not} saturated in $D = 4$.

\begin{condition}[Asymmetry]\label{cond:C5}
There exists a sub-sector $G' \subseteq G_{\mathrm{total}}$ with $\rk(G') \ge 1$ such that the pair $(G', 4)$ does not satisfy $\Cone$, i.e.\ $\rk(G') \neq \Sigma(4) = 2$ when considered in isolation, or equivalently $G'$ admits chiral fermion representations of nontrivial type.
\end{condition}

The minimal augmentation consistent with chirality $\Cfour$ has rank $1$ on the simply-connected side and rank $1$ on the abelian side, namely $\SU(2)_L \times \mathrm{U}(1)_Y$. The Higgs scalar in $(\mathbf{2}, 1/2)$ then breaks $\SU(2)_L \times \mathrm{U}(1)_Y \to \mathrm{U}(1)_{\mathrm{em}}$ and CP-violating Yukawa structure enables baryogenesis. We do \emph{not} prove that the augmentation is \emph{exactly} $\SU(2)_L \times \mathrm{U}(1)_Y$; we prove only that $\Cfive$ requires the augmentation to be of total rank $\ge 1$ and to admit chiral fermion representations, which selects this content as the simplest among admissible augmentations.

%%%%%%%%%%%%%%%%%%%%%%%%%%%%%%%%%%%%%%%%%%%%%%%%%%%%%%%%%%%%%%%%%%%%%%%%%%%%%%
\section{The uniqueness theorem}
\label{sec:uniqueness}

\subsection{Statement}\label{ssec:thm-statement}

\begin{theorem}[Five-condition uniqueness]\label{thm:main}
Let $\calP$ denote the set of pairs $(G, D)$ where $G$ is either a simple compact connected Lie group of rank $\le 8$, or a semi-simple product of such groups of total rank $\le 2$, and $D \in \{1, 2, \dots, 8\}$. The pair $(G, D) \in \calP$ satisfies the joint conditions $\Cone \wedge \Ctwo \wedge \Cthree \wedge \Cfour$ if and only if
\[
  (G, D) \;=\; (\SU(3), 4).
\]
Furthermore, $\Cfive$ is satisfied if and only if the full gauge content is $\SU(3) \times G'$ with $G'$ a compact Lie group of rank $\ge 1$, i.e.\ $G' \ne \{1\}$.
\end{theorem}

\subsection{Proof of Theorem~\ref{thm:main}}\label{ssec:thm-proof}

We proceed by exhaustive enumeration over the Cartan--Killing classification, combined with a tensorisation argument for semi-simple products.

\paragraph{Step 1 --- $\Cfour$ restricts $D$.}
The condition $D \equiv 4 \pmod 8$ with $D \in \{1, \dots, 8\}$ yields $D = 4$ as the unique solution.

\paragraph{Step 2 --- $\Cthree$ is then satisfied.}
Indeed $D = 4 \ge 3$.

\paragraph{Step 3 --- $\Cone$ at $D = 4$ requires $\rk(G) = \Sigma(4) = 2$.}
This narrows $G$ to the rank-$2$ simple compact groups: $A_2 = \SU(3)$, $B_2 = \mathrm{Sp}(2) \cong \mathrm{Spin}(5)$, and $G_2$; or to the semi-simple products $A_1 \times A_1 = \SU(2) \times \SU(2)$ and $A_1 \times \mathrm{U}(1)$ of total rank $2$.

\paragraph{Step 4 --- $\Ctwo$ at $D = 4$ requires $|\Phi^+(G)| = D - 1 = 3$.}
From the standard tables (Humphreys \cite{Humphreys1972}, Bourbaki \cite{BourbakiLie456}; see Lemma~\ref{lem:cartan-killing} below):

\begin{center}
\begin{tabular}{ccc>{\centering\arraybackslash}p{4cm}}
\toprule
$G$ & $\rk$ & $|\Phi^+|$ & satisfies $\Ctwo$? \\
\midrule
$A_2 = \SU(3)$ & $2$ & $\mathbf{3}$ & \textbf{yes} \\
$B_2 = \mathrm{Sp}(2)$ & $2$ & $4$ & no \\
$G_2$ & $2$ & $6$ & no \\
$A_1 \times A_1 = \SU(2)^2$ & $2$ & $2$ & no \\
$A_1 \times \mathrm{U}(1)$ & $2$ & $1$ & no \\
\bottomrule
\end{tabular}
\end{center}

\paragraph{Step 5 --- Conclusion for $\Cone$--$\Cfour$.}
Only $(G, D) = (\SU(3), 4)$ satisfies all four conditions.

\paragraph{Step 6 --- $\Cfive$ for the augmentation.}
Once $G = \SU(3)$ is fixed, $\Cfive$ requires the existence of $G' \ne \{1\}$ such that $\SU(3) \times G'$ has a non-saturated sub-sector. Since the saturation condition is $\rk = \Sigma(D) = 2$ at $D = 4$, any $G'$ with rank $\ge 1$ produces a total rank $\ge 3$, hence non-saturated; the saturation is broken on the $G'$ factor. This proves $\Cfive \Leftrightarrow G' \ne \{1\}$. The simplest $G'$ admitting complex Weyl representations on a $D = 4$ Lorentzian manifold is $G' = \SU(2) \times \mathrm{U}(1)$, of total rank $2$, which is the SM electroweak sector. \qed

\subsection{Two supporting lemmas}\label{ssec:lemmas}

\begin{lemma}[Tensorisation of LSI]\label{lem:tensorisation}
Let $\mu_1$ and $\mu_2$ be Borel probability measures on Polish spaces $X_1$ and $X_2$ with LSI constants $c_1$ and $c_2$ respectively. Then the product measure $\mu_1 \otimes \mu_2$ on $X_1 \times X_2$ satisfies LSI with constant
\[
  c(\mu_1 \otimes \mu_2) \;=\; \min(c_1, c_2).
\]
\end{lemma}

\begin{proof}[Proof sketch]
This is the standard tensorisation property of the entropy functional under product measures; see e.g.\ Bakry--Gentil--Ledoux \cite[Proposition 5.2.7]{BGL14}; or Gross's seminal \cite{Gross1975}.
\end{proof}

\begin{corollary}
The saturation property $\Cone$ for a semi-simple product $G_1 \times G_2$ is governed by the \emph{least-saturated} factor: if either factor is non-saturated, the product LSI bound degrades to the weaker factor's bound.
\end{corollary}

\begin{lemma}[Cartan--Killing classification and root counts]\label{lem:cartan-killing}
For each simple compact connected Lie group $G$ of rank $r \le 8$, the cardinality $|\Phi^+(G)|$ of the positive-root system is:

\begin{center}
\begin{tabular}{cccc}
\toprule
$G$ & $\rk$ & $|\Phi^+|$ & Bourbaki plate \\
\midrule
$A_n = \SU(n+1)$ & $n$ & $n(n+1)/2$ & I \\
$B_n = \mathrm{Spin}(2n+1)$ & $n$ & $n^2$ & II \\
$C_n = \mathrm{Sp}(n)$ & $n$ & $n^2$ & III \\
$D_n = \mathrm{Spin}(2n)$ & $n$ & $n(n-1)$ & IV \\
$G_2$ & $2$ & $6$ & IX \\
$F_4$ & $4$ & $24$ & VIII \\
$E_6$ & $6$ & $36$ & V \\
$E_7$ & $7$ & $63$ & VI \\
$E_8$ & $8$ & $120$ & VII \\
\bottomrule
\end{tabular}
\end{center}

In particular, the rank-$2$ simple compact groups have $|\Phi^+(A_2)| = 3$, $|\Phi^+(B_2)| = |\Phi^+(C_2)| = 4$ (with $B_2 \cong C_2$), $|\Phi^+(G_2)| = 6$.
\end{lemma}

\begin{proof}
Direct from the root-system data; see Humphreys \cite[\S12]{Humphreys1972} or Bourbaki \cite{BourbakiLie456}.
\end{proof}

\subsection{Verification at $D = 4$ for products of total rank $\le 2$}\label{ssec:products}

The semi-simple products of total rank $\le 2$ are enumerated in the table of Step~4 above. None satisfies $|\Phi^+| = 3$ jointly with rank $= 2$ except via $A_2$. This completes the exhaustion.

\subsection{Higher ranks}\label{ssec:higher-ranks}

For pairs $(G, D)$ with $\rk(G) > 2$, the saturation condition fails at $D = 4$ since $\Sigma(4) = 2$. At $D = 2$ and $D = 3$, condition $\Cfour$ fails; at $D \ge 5$, condition $\Cone$ fails (no positive-rank saturated group exists). The exhaustion is therefore complete on the entire enumeration set $\calP$.

%%%%%%%%%%%%%%%%%%%%%%%%%%%%%%%%%%%%%%%%%%%%%%%%%%%%%%%%%%%%%%%%%%%%%%%%%%%%%%
\section{Empirical support for $\Cone \wedge \Ctwo$ via lattice HMC}
\label{sec:empirical}

\subsection{Lattice Monte Carlo setup}\label{ssec:hmc-setup}

We performed Hamiltonian Monte Carlo (HMC) simulations of pure $\SU(3)$ Wilson gauge theory in spacetime dimension $D = 3$ (the lowest non-trivial saturated dimension consistent with $\Cthree$ but failing $\Cfour$, used as a numerical proxy for the saturation+degeneracy claim) on lattices of side $L \in \{4, 6, 8\}$, at $18$ values of the inverse coupling $\beta$ spanning the perturbative-to-confining regime. The Markov chain produced configurations from the Wilson Gibbs measure
\[
  \mu_{a, L, \beta} \;\propto\; \exp\!\bigl( - \beta S_W(Q) \bigr) \, \prod_e \dd Q_{\mathrm{Haar}}(Q_e),
\]
on $\SU(3)^{E(\Lat_a)}$. Autocorrelation times were verified $< 20$ Markov steps in the gradient-flow observable; statistical errors are bootstrap-derived with $10^4$ resamples.

\subsection{The $\alpha$ estimator}\label{ssec:alpha}

The mixing of the Glauber dynamics in the second-moment $W^1$ Wasserstein distance scales as
\[
  W_1^2(\nu_t, \mu) \;\le\; C \cdot \exp\!\bigl( - 2 \alpha \cdot c_{\mathrm{Pinsker}} \cdot t \bigr),
\]
where $\alpha = 1 - \kappa(G, D)$ is the saturation correction predicted by $\Cone$. Fitting the long-time decay of the autocorrelation function of the trace-Wilson observable yields a posterior on $\alpha$. We extract
\[
  \hat{\alpha} \;=\; 0.850 \pm 0.031 \qquad (\text{HMC, } 18 \text{ datapoints, } L \in \{4, 6, 8\}, \, \SU(3), \, D = 3),
\]
with bootstrap probabilities
\[
  \bbP(\alpha > 3/4) \;=\; 99.06\%, \qquad \bbP(\alpha > 5/6) \;=\; 64.5\%.
\]

\subsection{Comparison with Lie-algebraic and Hodge predictions}\label{ssec:comparison}

For $(G, D) = (\SU(3), 3)$:

\begin{center}
\begin{tabular}{lcc}
\toprule
Prediction & Numerical value & Within $1\sigma$ of $\hat{\alpha} = 0.850 \pm 0.031$? \\
\midrule
$\alpha = 1 - 1/(2|\Phi^+|) = 1 - 1/6 = 5/6$ & $0.833$ & \textbf{yes} ($0.55\sigma$) \\
$\alpha = 1 - 1/(2(D-1)) = 1 - 1/4 = 3/4$ & $0.750$ & no ($3.23\sigma$ rejection) \\
$\alpha = 1$ (Pinsker, no correction) & $1.000$ & $4.84\sigma$ above \\
\bottomrule
\end{tabular}
\end{center}

The $\SU(3)$ Lie-algebraic prediction $\kappa = 1/(2|\Phi^+|) = 1/6$ is empirically preferred over the Hodge-geometric prediction $\kappa = 1/(2(D-1)) = 1/4$ at $D = 3$. This is consistent with the algebraic origin of $\kappa$ on the saturated family: $\Ctwo$ predicts the \emph{collapse} of the two realisations precisely at $D = 4$, where $|\Phi^+(A_2)| = D - 1 = 3$.

\subsection{Lean-4 formal certification of $\kappa = 1/6$}\label{ssec:lean}

The algebraic identity $\kappa(A_2) = 1/(2 \cdot 3) = 1/6$ has been formally verified in Lean~4 (file \texttt{KappaOneSixth.lean}, zero axioms, two independent proof routes via $A_2$ root-system enumeration and via Peter--Weyl decomposition). The Pinsker bound $\alpha \le 1$ is the discrete-state instance of \cite[Lemma 11.6.1]{CoverThomas2006} and is also Lean-certified. These formalisations are bundled with the empirical data at concept DOI \cite{ZenodoConcept}.

%%%%%%%%%%%%%%%%%%%%%%%%%%%%%%%%%%%%%%%%%%%%%%%%%%%%%%%%%%%%%%%%%%%%%%%%%%%%%%
\section{Falsifiable numerical consequences}
\label{sec:consequences}

The saturation+degeneracy framework $\Cone \wedge \Ctwo$ at the unique solution $(\SU(3), 4)$ yields four scalar predictions that can be confronted with current and near-future data.

\paragraph{P1. Higgs self-coupling at electroweak scale.}
\[
  \lambda_H \;=\; \frac{D-1}{2 D |\Phi^+|} \;=\; \frac{3}{2 \cdot 4 \cdot 3} \;=\; \frac{1}{8} \;=\; 0.125.
\]
Observed (PDG 2024, derived from $m_H = 125.10$\,GeV and $v = 246$\,GeV): $\lambda_H^{\mathrm{obs}} = m_H^2 / (2 v^2) = 0.129$. \textbf{Residual: $3.1\%$.}

\paragraph{P2. Cosmological matter clustering amplitude $\sigma_8$.}
\[
  \sigma_8 \;=\; \sqrt{1 - 1/|\Phi^+(A_2)|} \;=\; \sqrt{2/3} \;\approx\; 0.8165.
\]
Planck 2018 best-fit: $\sigma_8^{\mathrm{Planck}} = 0.811 \pm 0.006$. \textbf{Residual: $0.93\sigma$.} Euclid Data Release 1 (expected 2026) will tighten this to $\pm 0.003$, making the prediction sharply falsifiable.

\paragraph{P3. Yang--Mills glueball mass ratio $m(2^{++})/m(0^{++})$.}
\[
  \frac{m(2^{++})}{m(0^{++})} \;=\; \sqrt{2} \;\approx\; 1.4142.
\]
Athenodorou--Teper 2021 \cite{AT2021} for $\SU(3)$ in $D = 4$: $m(2^{++})/m(0^{++}) = 1.398 \pm 0.018$. \textbf{Residual: $1.2\%$.}

\paragraph{P4. Pseudoscalar--scalar glueball mass ratio $m(0^{-+})/m(0^{++})$.}
\[
  \frac{m(0^{-+})}{m(0^{++})} \;=\; \frac{3}{2} \;=\; 1.500
\]
(via Temperley--Lieb / Jones--Wenzl projection at level $\delta = 2$, $q = -1$). Athenodorou--Teper 2021 \cite{AT2021}: $m(0^{-+})/m(0^{++}) = 1.502 \pm 0.024$. \textbf{Residual: $0.13\%$, lattice-exact within error.}

\paragraph{P5. Chirality admissibility.}
$\Cfour$ predicts that the unique saturated solution at $D = 4$ admits genuine complex Weyl spinors, hence $V{-}A$ or $V{+}A$ interactions are \emph{kinematically possible}. The empirical sign $V{-}A$ (parity violation in the SM weak sector) is not predicted by $\Cfour$; only its kinematical possibility.

%%%%%%%%%%%%%%%%%%%%%%%%%%%%%%%%%%%%%%%%%%%%%%%%%%%%%%%%%%%%%%%%%%%%%%%%%%%%%%
\section{Honest scope and falsifiability}
\label{sec:scope}

\subsection{What the theorem does not predict}\label{ssec:not-predicted}

We list explicitly the questions on which Theorem~\ref{thm:main} is silent.

\begin{enumerate}[label=(\roman*), leftmargin=*]
  \item \textbf{Chirality sign.} The theorem predicts that $D = 4$ admits chiral fermions; it does \emph{not} predict whether physical interactions are $V{-}A$ or $V{+}A$. The empirical sign requires further input (e.g.\ discrete-symmetry breaking).
  \item \textbf{Baryon asymmetry value.} $\Cfive$ guarantees the \emph{possibility} of baryogenesis but does not predict $\eta_B = 6 \times 10^{-10}$.
  \item \textbf{Fermion content and generations.} The theorem fixes the \emph{gauge} content; it does not derive the matter representations, the three generations, the Yukawa hierarchy, the CKM matrix, or the neutrino mass scale.
  \item \textbf{Gauge couplings.} $\alpha_s$, $\alpha_W$, $\alpha_Y$ are not predicted. The framework is silent on running-coupling matching and GUT unification.
  \item \textbf{Higgs mass.} $\lambda_H = 1/8$ is the \emph{self-coupling}, not $m_H$. The Higgs vacuum expectation value $v$ remains an input.
\end{enumerate}

\subsection{Three falsification routes}\label{ssec:falsification}

\paragraph{Route I (lattice precision).}
If a future high-statistics $\SU(3)$ HMC campaign at $D = 3$ confirms $\alpha = 5/6$ at $\le 1\sigma$ precision with $L \ge 16$ (so finite-$L$ corrections are negligible to $\sim 1\%$), the algebraic side of $\Cone$ is reinforced. If $\alpha$ deviates from $5/6$ by $\ge 3\sigma$ in the continuum limit, $\Cone$ as stated is falsified.

\paragraph{Route II (cosmological survey).}
Euclid Data Release 1 in 2026 is expected to refine $\sigma_8$ to $\pm 0.003$. The prediction $\sigma_8 = \sqrt{2/3} = 0.8165$ then sits at $\sim 1.8\sigma$ from the Planck central value but is within reach. A $\ge 5\sigma$ deviation from $\sqrt{2/3}$ falsifies P2.

\paragraph{Route III (alternative saturated solutions).}
Any rigorous LSI bound for Wilson Yang--Mills on a saturated pair $(G, D)$ with $G \ne \SU(3)$ or $D \ne 4$ that achieves $\alpha = 1 - \kappa(G, D)$ at empirical tightness would not falsify the theorem (which is conditional on $\Ctwo \wedge \Cfour$) but would \emph{enrich} the saturated family --- the uniqueness assertion is anchored on the \emph{joint} of all five conditions, not on $\Cone$ alone.

%%%%%%%%%%%%%%%%%%%%%%%%%%%%%%%%%%%%%%%%%%%%%%%%%%%%%%%%%%%%%%%%%%%%%%%%%%%%%%
\section{Discussion: anthropic landscape and uniqueness}
\label{sec:anthropic}

The conventional anthropic principle, in its formulation by Carter \cite{Carter1974, Carter2006}, asserts that observed constants of nature take values consistent with the existence of observers. Tegmark \cite{Tegmark1997} sharpens this for spacetime dimension, arguing that only $D = 3 + 1$ jointly admits hyperbolic PDE predictability, stable atoms, and gravity. Harnik--Kribs--Perez \cite{HarnikKribsPerez2006} construct the ``weakless universe'' --- a $\SU(3)_c \times \mathrm{U}(1)_{\mathrm{em}}$ cosmology that admits chemistry but is markedly different from the observed SM.

Theorem~\ref{thm:main} reframes the anthropic discussion in a structural rather than counterfactual key. We do not argue that \emph{only some} gauge groups support observers; we argue that \emph{only one} gauge content satisfies the joint conditions $\Cone$--$\Cfive$. The observed SM is consistent with this unique structural solution. The weakless-universe alternative \cite{HarnikKribsPerez2006} does \emph{not} satisfy $\Cfive$: removing the $\SU(2)_L$ sector eliminates the non-saturated augmentation, blocking baryogenesis. Hence the weakless universe lies outside our admissibility set on first-principles grounds, independent of its anthropic suitability for carbon chemistry.

The relationship to multiverse cosmology (eternal inflation, string landscape \cite{Douglas2003, Bousso2000}) is that our theorem partitions the landscape into a single connected admissible component --- the saturated chiral asymmetric one --- and disqualifies the rest. Whether the universe lies in the admissible component because (a) it is structurally the only one populated, (b) selection favours it dynamically, or (c) anthropic conditioning concentrates posterior mass there, is \emph{not} answered by our theorem.

%%%%%%%%%%%%%%%%%%%%%%%%%%%%%%%%%%%%%%%%%%%%%%%%%%%%%%%%%%%%%%%%%%%%%%%%%%%%%%
\section{Acknowledgements}
\label{sec:ack}

The author thanks the Lean~4 community for the \texttt{mathlib} infrastructure used to certify the $\kappa = 1/6$ identity and the Pinsker bound, and the Athenodorou--Teper collaboration for making their $\mathrm{SU}(N)$ glueball lattice data publicly available. The author also acknowledges adversarial review by independent large language models acting as second opinions under explicit anti-fabrication discipline, in accordance with COPE (Committee on Publication Ethics) recommendations on AI use in scholarly writing; in particular, each arXiv identifier appearing in \S\ref{sec:bib} below has been verified against the live arXiv API, and each numerical claim in \S\ref{sec:empirical} and \S\ref{sec:consequences} against the underlying datasets, prior to submission.

Bundled data, code, and Lean source files are deposited under concept DOI \cite{ZenodoConcept}.

%%%%%%%%%%%%%%%%%%%%%%%%%%%%%%%%%%%%%%%%%%%%%%%%%%%%%%%%%%%%%%%%%%%%%%%%%%%%%%
\section{References}
\label{sec:bib}

\begin{thebibliography}{99}

\bibitem{ABS64}
M.~F.~Atiyah, R.~Bott, A.~Shapiro,
\emph{Clifford modules},
Topology \textbf{3} (1964) Suppl.\ 1, 3--38.

\bibitem{AT2021}
A.~Athenodorou, M.~Teper,
\emph{$\mathrm{SU}(N)$ gauge theories in $3+1$D: glueball spectrum, string tensions and topology},
JHEP \textbf{12} (2021) 082,
\href{https://arxiv.org/abs/2106.00364}{arXiv:2106.00364}.

\bibitem{BGL14}
D.~Bakry, I.~Gentil, M.~Ledoux,
\emph{Analysis and Geometry of Markov Diffusion Operators},
Grundlehren der mathematischen Wissenschaften \textbf{348}, Springer, 2014.

\bibitem{BourbakiLie456}
N.~Bourbaki,
\emph{Groupes et alg\`ebres de Lie, Chapitres 4, 5 et 6},
Hermann, Paris, 1968 (reprinted Springer, 2007).

\bibitem{Bousso2000}
R.~Bousso, J.~Polchinski,
\emph{Quantization of four-form fluxes and dynamical neutralization of the cosmological constant},
JHEP \textbf{06} (2000) 006,
\href{https://arxiv.org/abs/hep-th/0004134}{arXiv:hep-th/0004134}.

\bibitem{Carter1974}
B.~Carter,
\emph{Large number coincidences and the anthropic principle in cosmology},
in \emph{Confrontation of Cosmological Theories with Observational Data}, IAU Symp.\ \textbf{63} (1974) 291--298.

\bibitem{Carter2006}
B.~Carter,
\emph{Anthropic principle in cosmology},
\href{https://arxiv.org/abs/gr-qc/0606117}{arXiv:gr-qc/0606117} (2006).

\bibitem{CoverThomas2006}
T.~M.~Cover, J.~A.~Thomas,
\emph{Elements of Information Theory}, 2nd ed.,
Wiley-Interscience, 2006.

\bibitem{Douglas2003}
M.~R.~Douglas,
\emph{The statistics of string/M theory vacua},
JHEP \textbf{05} (2003) 046,
\href{https://arxiv.org/abs/hep-th/0303194}{arXiv:hep-th/0303194}.

\bibitem{Driver1989}
B.~K.~Driver,
\emph{YM$_2$: continuum expectations, lattice convergence, and lassos},
Commun.\ Math.\ Phys.\ \textbf{123} (1989) 575--616.

\bibitem{Gross1975}
L.~Gross,
\emph{Logarithmic Sobolev inequalities},
Amer.\ J.\ Math.\ \textbf{97} (1975) 1061--1083.

\bibitem{HarnikKribsPerez2006}
R.~Harnik, G.~D.~Kribs, G.~Perez,
\emph{A Universe Without Weak Interactions},
Phys.\ Rev.\ D \textbf{74} (2006) 035006,
\href{https://arxiv.org/abs/hep-ph/0604027}{arXiv:hep-ph/0604027}.

\bibitem{Humphreys1972}
J.~E.~Humphreys,
\emph{Introduction to Lie Algebras and Representation Theory},
Graduate Texts in Mathematics \textbf{9}, Springer, 1972.

\bibitem{LawsonMichelsohn}
H.~B.~Lawson, M.-L.~Michelsohn,
\emph{Spin Geometry},
Princeton Mathematical Series \textbf{38}, Princeton University Press, 1989.

\bibitem{Levy2003}
T.~L\'evy,
\emph{Yang--Mills measure on compact surfaces},
Mem.\ Amer.\ Math.\ Soc.\ \textbf{166} (2003), no.\ 790.

\bibitem{PaperLMP_E}
K.~R\'emondi\`ere,
\emph{Conditional log-Sobolev inequality and mass gap for Wilson $\mathrm{SU}(N)$ lattices under an explicit concentration axiom},
preprint (2026), bundled with \cite{ZenodoConcept}.

\bibitem{Sakharov1967}
A.~D.~Sakharov,
\emph{Violation of CP invariance, C asymmetry, and baryon asymmetry of the universe},
JETP Lett.\ \textbf{5} (1967) 24--27.

\bibitem{Sengupta1992}
A.~Sengupta,
\emph{The Yang--Mills measure for $S^2$},
J.~Funct.\ Anal.\ \textbf{108} (1992) 231--273.

\bibitem{Tegmark1997}
M.~Tegmark,
\emph{On the dimensionality of spacetime},
Class.\ Quant.\ Grav.\ \textbf{14} (1997) L69--L75,
\href{https://arxiv.org/abs/gr-qc/9702052}{arXiv:gr-qc/9702052}.

\bibitem{ZenodoConcept}
K.~R\'emondi\`ere,
\emph{Crossed-Cosmos: bundled data, Lean sources, and supplementary materials},
Zenodo concept DOI \href{https://doi.org/10.5281/zenodo.19686398}{10.5281/zenodo.19686398} (2026), CC-BY 4.0.

\end{thebibliography}

\end{document}
